\section{Certificates for the deterministic lower bounds}
\label{app:deterministic-lower-bounds}

\subsection{The logarithmic-spiral construction}
\label{app:spiral}

This subsection proves \Cref{thm:global:spiral-lower}.  The construction is
two-dimensional; the extension to higher dimension is recorded at the end.

\subsubsection{Potential and global Hessian certificate}

Fix \(K\geq16\) and set
\begin{equation}
\label{app:spiral:parameters}
    \alpha_0:=2\alpha,
    \qquad
    c:=\frac{\sqrt{K-2}}3,
    \qquad
    \delta_K:=\frac{10^{-4}}K,
    \qquad
    D_K:=\frac{18\alpha_0(K-1)}{K+7},
    \qquad
    \beta:=\alpha(2K+1).
\end{equation}
Thus the condition number in the original coordinates is \(2K+1\).
Define
\[
    s(u):=35u^4-84u^5+70u^6-20u^7,
    \qquad 0\leq u\leq1,
\]
and the even \(C^3\) cutoff
\begin{equation}
\label{app:spiral:phase-bump}
    \rho(u):=
    \begin{cases}
        1, & |u|\leq\tfrac12,\\[1mm]
        s(2-2|u|), & \tfrac12<|u|<1,\\[1mm]
        0, & |u|\geq1.
    \end{cases}
\end{equation}
Extend \(\rho_{\delta_K}(\phi):=\rho(\phi/\delta_K)\) periodically with
period \(\pi\), and put
\[
    \bar\rho_{\delta_K}
    :=\frac1\pi\int_{-\pi/2}^{\pi/2}\rho_{\delta_K}(\phi)\,\dd\phi,
    \qquad
    g_K(\phi)
    :=\frac{\rho_{\delta_K}(\phi)-\bar\rho_{\delta_K}}
            {1-\bar\rho_{\delta_K}}.
\]
Let \(f_K\) be the even \(\pi\)-periodic solution of
\begin{equation}
\label{app:spiral:fK-equation}
    f_K''=g_K,
    \qquad
    f_K(0)=f_K'(0)=0.
\end{equation}
The normalization gives
\begin{equation}
\label{app:spiral:fK-bounds}
    f_K(\phi)=\frac12\phi^2
    \quad\text{for }|\phi|\leq\frac{\delta_K}{2},
    \qquad
    \norm{f_K'}_\infty\leq3\delta_K,
    \qquad
    \norm{f_K}_\infty\leq2\pi\delta_K.
\end{equation}

For the radial cutoff use
\begin{equation}
\label{app:spiral:radial-cutoff}
    \chi(\ell):=
    \begin{cases}
        0, & \ell\leq0,\\
        35\ell^4-84\ell^5+70\ell^6-20\ell^7,
            & 0<\ell<1,\\
        1, & \ell\geq1.
    \end{cases}
\end{equation}
Fix a scale \(r_0>0\).  In polar coordinates \((r,\theta)\), set
\[
    \ell:=\log\frac r{r_0},
    \qquad
    \phi:=\theta+c\ell,
\]
and define
\begin{equation}
\label{app:spiral:potential}
    \boxed{
    V_K(r,\theta)
    :=\frac12r^2\bigl[\alpha_0+D_K\chi(\ell)f_K(\phi)\bigr].}
\end{equation}
For \(r\leq r_0\), the potential is exactly
\(V_K(x)=\alpha_0\norm{x}^2/2\).  The first three derivatives of \(\rho\) and
\(\chi\) match at the endpoints of their pieces, so
\(V_K\in C^3(\R^2)\).  Its evenness follows from the \(\pi\)-periodicity of
\(f_K\).

\begin{lemma}[Global Hessian certificate]
\label{app:spiral:hessian-certificate}
The potential in \eqref{app:spiral:potential} satisfies
\begin{equation}
\label{app:spiral:hessian-bounds}
    0.98\alpha_0\Id\preceq\nabla^2V_K(x)\preceq \alpha_0(K+0.02)\Id
    \qquad (x\in\R^2).
\end{equation}
Consequently,
\(\alpha\Id\preceq\nabla^2V_K\preceq\beta\Id\), with \(\beta\) from
\eqref{app:spiral:parameters}.  On the flat part of a spiral valley, the
two Hessian eigenvalues are exactly \(\alpha_0\) and \(\alpha_0K\).
\end{lemma}

\begin{proof}
Put \(q(\ell,\theta):=\chi(\ell)f_K(\theta+c\ell)\).  For any profile
\(A=A(\ell,\theta)\), direct differentiation in the orthonormal polar
frame yields
\[
    \nabla^2\left(\frac12r^2A\right)
    =
    \begin{pmatrix}
      A+\frac32A_\ell+\frac12A_{\ell\ell}
      &\frac12A_\theta+\frac12A_{\ell\theta}\\[1mm]
      \frac12A_\theta+\frac12A_{\ell\theta}
      &A+\frac12A_\ell+\frac12A_{\theta\theta}
    \end{pmatrix}.
\]
It follows that
\begin{equation}
\label{app:spiral:hessian-decomposition}
 \nabla^2V_K-\alpha_0\Id
 =\frac{D_K}{2}\chi g_K
     \begin{pmatrix}c^2&c\\c&1\end{pmatrix}
   +D_K\{\chi B_K+E_K\},
\end{equation}
where all scalar functions are evaluated at \((\ell,\phi)\) and
\[
 B_K=
 \begin{pmatrix}
 f_K+\frac32cf_K'&\frac12f_K'\\
 \frac12f_K'&f_K+\frac12cf_K'
 \end{pmatrix},
 \qquad
 E_K=
 \begin{pmatrix}
 (\frac32\chi'+\frac12\chi'')f_K+c\chi'f_K'&\frac12\chi'f_K'\\
 \frac12\chi'f_K'&\frac12\chi'f_K
 \end{pmatrix}.
\]
This identity includes the transition annulus \(0<\ell<1\).

The normalization of the phase cutoff is explicit.  Since
\(\int_{-1}^1\rho(u)\,\dd u=3/2\),
\[
 \bar\rho_{\delta_K}=\frac{3\delta_K}{2\pi},
 \qquad
 -\frac{\bar\rho_{\delta_K}}{1-\bar\rho_{\delta_K}}
 \leq g_K\leq1.
\]
For \(K\geq16\),
\begin{equation}
\label{app:spiral:negative-background}
 (K-1)\frac{\bar\rho_{\delta_K}}{1-\bar\rho_{\delta_K}}
 <4.8\cdot10^{-5}.
\end{equation}
The radial cutoff satisfies
\[
 0\leq\chi\leq1,
 \qquad
 \norm{\chi'}_\infty=\frac{35}{16},
 \qquad
 \norm{\chi''}_\infty<7.514.
\]
Using \eqref{app:spiral:fK-bounds}, \(D_K/\alpha_0\leq18\), and
\(c\leq\sqrt K/3\), the maximum-row-sum bound in
\eqref{app:spiral:hessian-decomposition} gives
\begin{align*}
 \frac{D_K}{\alpha_0}\norm{\chi B_K+E_K}_{\op}
 &\leq18\Bigl[
   \left(1+\frac32\frac{35}{16}+\frac{7.514}{2}\right)2\pi\delta_K\\
 &\hspace{25mm}
  +\left(\frac32c+\frac12+c\frac{35}{16}
         +\frac12\frac{35}{16}\right)3\delta_K
  \Bigr]
 <0.0078.
\end{align*}
Every term in the last bound decreases with \(K\), so evaluation at \(K=16\)
suffices.

The rank-one matrix in the first term of
\eqref{app:spiral:hessian-decomposition} has nonzero eigenvalue
\(1+c^2=(K+7)/9\), and
\[
    \frac{D_K}{2}(1+c^2)=\alpha_0(K-1).
\]
Its positive part is at most \(\alpha_0(K-1)\Id\), while
\eqref{app:spiral:negative-background} bounds its negative part by
\(4.8\cdot10^{-5}\alpha_0\Id\).  Thus the calculation in fact proves
\[
 0.992\alpha_0\Id\preceq\nabla^2V_K
 \preceq \alpha_0(K+0.0078)\Id,
\]
which implies \eqref{app:spiral:hessian-bounds}.  On an uncut valley core,
\(g_K=1\) and \(f_K(0)=f_K'(0)=0\), so the two eigenvalues are \(\alpha_0\) and
\(\alpha_0K\).  Finally, \(0.98\alpha_0=1.96\alpha>\alpha\) and
\(\alpha_0(K+0.02)<\alpha(2K+1)=\beta\).
\end{proof}

\subsubsection{Reference orbit}

Let \(R_\theta\) be rotation through angle \(\theta\), put
\(s_K:=\sqrt{K-2}\), and define
\begin{equation}
\label{app:spiral:QK}
    Q_K
    :=\frac1{K+7}
      \begin{pmatrix}
        (K+1)^2 & 2s_K(K+1)\\
        2s_K(K+1) & 6(K+1)
      \end{pmatrix}.
\end{equation}
Its trace and determinant are
\[
    \Tr Q_K=K+1,
    \qquad
    \det Q_K=\frac{2(K+1)^2}{K+7},
\]
and
\[
 \det\left(Q_K-\frac32\Id\right)
 =\frac{2K^2-23K+29}{4(K+7)}>0
 \qquad(K\geq16).
\]
The principal minors of \(Q_K-3\Id/2\) and \(K\Id-Q_K\) therefore give
\begin{equation}
\label{app:spiral:QK-band}
    \frac32\Id\prec Q_K\prec K\Id,
    \qquad
    Q_K^{-1}e_1
    =\frac1{K+1}\begin{pmatrix}3\\-\sqrt{K-2}\end{pmatrix}.
\end{equation}

Choose \(\norm{m_0}=e^2r_0\) and an angle \(\theta_0\) such that
\begin{equation}
\label{app:spiral:phase-initialization}
    \theta_0+c\log\frac{\norm{m_0}}{r_0}\equiv0\pmod\pi,
\end{equation}
and initialize the precision by
\begin{equation}
\label{app:spiral:precision-initialization}
    P_0:=C_0^{-1}
    :=\alpha_0R_{\theta_0}Q_KR_{\theta_0}^{\top}.
\end{equation}
By \eqref{app:spiral:QK-band},
\begin{equation}
\label{app:spiral:initial-covariance-band}
    \frac1\beta\Id\preceq C_0\preceq\frac1\alpha\Id.
\end{equation}

The point-system reference orbit is
\begin{equation}
\label{app:spiral:reference-orbit}
    r_t^{\mathrm{ref}}
    :=e^2r_0\exp\left(-\frac{3t}{K+1}\right),
    \qquad
    \theta_t^{\mathrm{ref}}
    :=\theta_0+\frac{\sqrt{K-2}}{K+1}t,
\end{equation}
with
\[
    P_t^{\mathrm{ref}}
    :=\alpha_0R_{\theta_t^{\mathrm{ref}}}Q_K
      R_{\theta_t^{\mathrm{ref}}}^{\top}.
\]
Indeed, on the valley \(\nabla V_K(m)=hm\), and
\eqref{app:spiral:QK-band} gives
\[
    \dot r=-\frac3{K+1}r,
    \qquad
    r\dot\theta=\frac{\sqrt{K-2}}{K+1}r.
\]
A direct substitution also verifies
\(\dot P^{\mathrm{ref}}
=\nabla^2V_K(m^{\mathrm{ref}})-P^{\mathrm{ref}}\).

\subsubsection{Gaussian shadowing and the first-exit argument}

\begin{lemma}[Gaussian shadowing certificate]
\label{app:spiral:shadowing}
Assume
\begin{equation}
\label{app:spiral:radius-condition}
    r_0\sqrt\alpha
    \geq10^{25}K^{3/2}\sqrt{\log(eK)}.
\end{equation}
Let \((m_t,C_t)\) be the exact Gaussian Fisher--Rao flow for \(V_K\),
initialized by
\eqref{app:spiral:phase-initialization}--%
\eqref{app:spiral:precision-initialization}.  Then, for
\(0\leq t\leq K/100\),
\begin{equation}
\label{app:spiral:radial-shadowing}
    -\frac4K
    \leq\frac{\dd}{\dd t}\log\norm{m_t}
    \leq-\frac2K.
\end{equation}
In particular,
\begin{equation}
\label{app:spiral:mean-lower}
    \norm{m_t}\geq e^{-0.04}\norm{m_0}
    \qquad(0\leq t\leq K/100).
\end{equation}
\end{lemma}

\begin{proof}
Write \(P_t=C_t^{-1}\), choose the continuous lift of the polar angle
starting at \(\theta_0\), and introduce
\begin{equation}
\label{app:spiral:transverse-coordinates}
 z_t:=\theta_t+c\log\frac{r_t}{r_0},
 \qquad
 S_t:=\frac1{\alpha_0}R_{\theta_t}^{\top}P_tR_{\theta_t},
 \qquad
 E_t:=S_t-Q_K.
\end{equation}
The phase \(z_t\) is represented near zero modulo \(\pi\).  Put
\[
 w:=\begin{pmatrix}c\\1\end{pmatrix},
 \qquad
 J:=\begin{pmatrix}0&-1\\1&0\end{pmatrix},
 \qquad
 d_K:=\frac{D_K}{\alpha_0}=\frac{18(K-1)}{K+7}.
\]
On the uncut flat phase core \(f_K(z)=z^2/2\).  The normalized point
gradient and Hessian there are
\begin{align}
\label{app:spiral:core-point-gradient}
 b_K(z)
 &:=\frac1{\alpha_0 r}R_\theta^\top\nabla V_K(m)
 =\begin{pmatrix}
 1+\frac{d_K}{2}z^2+\frac{d_Kc}{2}z\\[1mm]
 \frac{d_K}{2}z
 \end{pmatrix},\\
\label{app:spiral:core-point-hessian}
 \mathsf H_K(z)
 &:=\frac1{\alpha_0}R_\theta^\top\nabla^2V_K(m)R_\theta\notag\\
 &=\Id+\frac{d_K}{2}ww^\top
 +d_K\begin{pmatrix}
 \frac12z^2+\frac32cz&\frac12z\\[1mm]
 \frac12z&\frac12z^2+\frac12cz
 \end{pmatrix}.
\end{align}
For the exact Gaussian flow define
\[
 \bar b_t:=\frac1{\alpha_0 r_t}R_{\theta_t}^\top
            \E[\nabla V_K(X_t)],
 \qquad
 \bar H_t:=\frac1{\alpha_0}R_{\theta_t}^\top
            \E[\nabla^2V_K(X_t)]R_{\theta_t},
 \qquad
 q_t:=S_t^{-1}\bar b_t.
\]
The mean and precision equations become
\begin{equation}
\label{app:spiral:rotated-system}
 \frac{\dot r_t}{r_t}=-(q_t)_1,
 \qquad
 \dot\theta_t=-(q_t)_2,
 \qquad
 \dot z_t=-w^\top q_t,
 \qquad
 \dot S_t=\bar H_t-S_t-(q_t)_2(S_tJ-JS_t).
\end{equation}
If the Gaussian averages are replaced by
\eqref{app:spiral:core-point-gradient}--%
\eqref{app:spiral:core-point-hessian}, then
\((z,S)=(0,Q_K)\) is an equilibrium.  With
\(\omega_K=s_K/(K+1)\),
\[
 Q_K^{-1}e_1
 =\begin{pmatrix}3/(K+1)\\-\omega_K\end{pmatrix},
 \qquad
 w^\top Q_K^{-1}e_1=0,
 \qquad
 \mathsf H_K(0)-Q_K
 +\omega_K(Q_KJ-JQ_K)=0.
\]

We next certify the linearization and nonlinear remainder.  Let
\begin{equation}
\label{app:spiral:y-scaling}
 \mathsf T_K:=\diag(K,\sqrt K,1,\sqrt K),
 \qquad
 y:=\mathsf T_K(z,E_{11},E_{12},E_{22})^\top.
\end{equation}
For \(S=Q_K+E\), let \(\mathcal F_K(y)\) be the scaled point vector field
obtained from \eqref{app:spiral:rotated-system} with
\(\bar b=b_K(z)\) and \(\bar H=\mathsf H_K(z)\).  Set
\[
 A_K:=D\mathcal F_K(0),
 \qquad
 \mathcal N_K(y):=\mathcal F_K(y)-A_Ky.
\]
Writing \(s=\sqrt{K-2}\), direct differentiation gives
\begin{equation}
\label{app:spiral:linear-matrix}
 A_K=\mathsf T_K\mathsf B_K\mathsf T_K^{-1},
\end{equation}
where, in the order \((z,E_{11},E_{12},E_{22})\),
{\scriptsize
\[
\mathsf B_K=
\begin{pmatrix}
-\frac{3(K-1)}{2(K+1)}&0&\frac{K+7}{2(K+1)^2}
&-\frac{s(K+7)}{6(K+1)^2}\\[1mm]
\frac{3s(K-1)}{K+7}
&-\frac{K^2+20K-17}{(K+1)(K+7)}
&\frac{12s}{K+7}&-\frac{2(K-2)}{K+7}\\[1mm]
\frac{3(K-1)(K+1)}{2(K+7)}
&\frac{2(K-11)s}{(K+1)(K+7)}
&-\frac{7K^2-10K+19}{2(K+1)(K+7)}
&\frac{s(K^2-2K+9)}{2(K+1)(K+7)}\\[1mm]
\frac{9s(K-1)}{K+7}
&\frac{12(K-2)}{(K+1)(K+7)}
&-\frac{12s}{K+7}&\frac{K-11}{K+7}
\end{pmatrix}.
\]
}
Its characteristic polynomial factors as
\begin{equation}
\label{app:spiral:characteristic-polynomial}
 \det(\lambda\Id-A_K)
 =\frac{\lambda+1}{2(K+1)^3}\,p_K(\lambda+1),
\end{equation}
where
\[
 p_K(\mu)
 =2(K+1)^3\mu^3
 +2(K-5)(K+1)^2\mu^2
 +4(K-2)(K+1)^2\mu
 +(K-2)(K^2-10K-23).
\]
For \(K\geq16\), all coefficients are positive and the cubic
Routh--Hurwitz determinant is
\[
 6(K-2)(K-1)^2(K+1)^3>0.
\]
Thus every eigenvalue of \(A_K\) has real part at most \(-1\).  The
entries of the scaled matrix satisfy
\[
 \bigl(|(A_K)_{ij}|\bigr)
 \preceq
 \begin{pmatrix}
 3/2&0&2/3&1/4\\
 3&3/2&12&2\\
 3/2&1/8&7/2&1/2\\
 9&3/4&12&1
 \end{pmatrix},
\]
so \(\norm{A_K}_{\op}<21\).  In a complex Schur decomposition
\(A_K=U(D+N)U^*\), the matrix \(N\) is strictly upper triangular,
\(\norm{N}_{\op}\leq42\), and
\(\norm{e^{Dt}}_{\op}\leq e^{-t}\).  Iterated Duhamel expansion
terminates after three factors of \(N\) and yields
\begin{equation}
\label{app:spiral:semigroup-bound}
 \norm{e^{A_Kt}}_{\op}
 \leq e^{-t}\sum_{j=0}^3\frac{(42t)^j}{j!}
 <2\cdot10^6e^{-t/2}.
\end{equation}

The nonlinear estimate is a finite-dimensional calculation.  Insert the
explicit inverse of \(S\) into
\eqref{app:spiral:core-point-gradient}--%
\eqref{app:spiral:core-point-hessian}.  For
\(\norm{y}\leq10^{-12}\) and \(K\geq16\),
\begin{equation}
\label{app:spiral:second-derivative-table}
 \sup\sum_{j,\ell=1}^4
 \left|\partial_{j\ell}^2(\mathcal F_K)_i\right|
 \leq(100,2000,500,2000)_i.
\end{equation}
Indeed, these bounds follow term by term from
\(d_K\leq18\), \(c/\sqrt K\leq1/3\),
\(Q_K\succeq3\Id/2\), and
\[
 E=\begin{pmatrix}
 y_2/\sqrt K&y_3\\y_3&y_4/\sqrt K
 \end{pmatrix};
\]
on this ball,
\(\norm{(Q_K+E)^{-1}}_{\op}\leq1\).  Taylor's theorem then gives
\begin{equation}
\label{app:spiral:nonlinear-bound}
 \norm{\mathcal N_K(y)}
 \leq\frac12\sqrt{100^2+2000^2+500^2+2000^2}\,\norm{y}^2
 <4000\norm{y}^2.
\end{equation}

It remains to compare Gaussian averages with the point fields at the actual
state.  Define the first-exit time
\begin{equation}
\label{app:spiral:stopping-time}
 \tau:=\frac K{100}\wedge
 \inf\left\{t\geq0:\norm{y_t}=10^{-12}
 \ \text{or}\ r_t\notin[e^{3/2}r_0,e^{5/2}r_0]\right\}.
\end{equation}
On \([0,\tau]\), \eqref{app:spiral:QK-band} and
\eqref{app:spiral:y-scaling} imply \(S_t\succeq\Id\), hence
\(C_t\preceq \alpha_0^{-1}\Id\).  Put \(Z_t=X_t-m_t\) and
\[
 \mathcal G_t
 :=\left\{\norm{Z_t}\leq\frac{\delta_Kr_t}{8\sqrt K}\right\},
 \qquad
 p_K:=\sup_{t\leq\tau}\Prob(\mathcal G_t^c),
 \qquad
 M_K:=\sup_{t\leq\tau}
 \E\left[\sqrt{\alpha_0}\norm{Z_t}\one_{\mathcal G_t^c}\right].
\]
On \(\mathcal G_t\), every point of the segment
\(m_t+uZ_t\), \(0\leq u\leq1\), satisfies
\[
 |\Delta\theta|+c|\Delta\ell|
 \leq2(1+c)\frac{\norm{Z_t}}{r_t}
 <\frac{\delta_K}{4}.
\]
Because
\(|z_t|\leq10^{-12}/K<\delta_K/4\) and
\(r_t\geq e^{3/2}r_0\), the whole segment remains in the uncut flat
phase core.  Under a common Gaussian representation,
\(\sqrt{\alpha_0}\norm{Z_t}\leq\norm{G}\) for
\(G\sim\N(0,\Id_2)\).  The exact two-dimensional radial tail gives
\begin{equation}
\label{app:spiral:tail-certificate}
 p_K
 \leq\exp\left(-\frac{\delta_K^2\alpha_0r_0^2}{128K}\right)
 \leq10^{-50}K^{-30},
 \qquad
 M_K\leq\sqrt{2p_K}\leq\sqrt2\,10^{-25}K^{-15}.
\end{equation}

On the flat core, differentiation of
\eqref{app:spiral:core-point-hessian} in Cartesian coordinates gives
\begin{equation}
\label{app:spiral:local-hessian-lipschitz}
 \operatorname{Lip}_{\op}(\nabla^2V_K)
 \leq\frac{15\alpha_0K}{r_t}
\end{equation}
along these segments.  For example,
\(\norm{\partial_z\mathsf H_K}_{\op}\leq12\sqrt K\),
\(\norm{\nabla z}\leq2\sqrt K/(5r)\), and rotation of the polar frame
contributes at most \(2\alpha_0(K+0.02)/r\); allowing a factor two for the
radius along the segment gives the displayed constant.

Define the normalized Gaussian errors
\begin{align*}
 \varepsilon_b(t)
 &:=\left\|\frac1{\alpha_0 r_t}R_{\theta_t}^\top
   \bigl(\E[\nabla V_K(X_t)]-\nabla V_K(m_t)\bigr)\right\|,\\
 \varepsilon_H(t)
 &:=\left\|\frac1{\alpha_0}R_{\theta_t}^\top
   \bigl(\E[\nabla^2V_K(X_t)]-\nabla^2V_K(m_t)\bigr)
   R_{\theta_t}\right\|_{\op}.
\end{align*}
Let \(A_0:=\sqrt\alpha\,r_0\).  Split the expectations at
\(\mathcal G_t\), use
\eqref{app:spiral:local-hessian-lipschitz} for the Taylor remainder on
\(\mathcal G_t\), the global Hessian bound on its complement,
\(\E Z_t=0\), and \(\E\norm{Z_t}^2\leq2/\alpha_0\).  This gives
\begin{equation}
\label{app:spiral:average-errors}
 \varepsilon_b(t)
 \leq8\frac K{A_0^2}+3\frac K{A_0}M_K,
 \qquad
 \varepsilon_H(t)
 \leq15\frac K{A_0}+3Kp_K.
\end{equation}
Let \(\xi_K(t)\) be the contribution of these actual-state errors to
the scaled transverse system.  Since
\(\norm{S_t^{-1}}_{\op}\leq1\),
\(\norm{S_t}_{\op}\leq2K\), and
\(\norm{w}\leq2\sqrt K/5\), equations
\eqref{app:spiral:rotated-system} and
\eqref{app:spiral:y-scaling} imply
\begin{align}
\label{app:spiral:smoothing-certificate}
 \norm{\xi_K(t)}
 &\leq\frac25K^{3/2}\varepsilon_b(t)
   +\sqrt{2K}\bigl(\varepsilon_H(t)+4K\varepsilon_b(t)\bigr)\notag\\
 &\leq23\frac{K^{3/2}}{A_0}
   +5K^{3/2}p_K
   +52\frac{K^{5/2}}{A_0^2}
   +20\frac{K^{5/2}}{A_0}M_K.
\end{align}
The radius condition \eqref{app:spiral:radius-condition} and
\eqref{app:spiral:tail-certificate} make the last line smaller than
\(3\cdot10^{-24}\), hence smaller than \(10^{-22}\), throughout
\([0,\tau]\).

The exact transverse equation is therefore
\begin{equation}
\label{app:spiral:transverse-system}
 \dot y_t=A_Ky_t+\mathcal N_K(y_t)+\xi_K(t),
 \qquad
 y_0=0.
\end{equation}
Variation of constants together with
\eqref{app:spiral:semigroup-bound},
\eqref{app:spiral:nonlinear-bound}, and
\eqref{app:spiral:smoothing-certificate} yields, for \(t\leq\tau\),
\[
 \norm{y_t}
 \leq2\cdot10^6\int_0^te^{-(t-s)/2}
 \bigl(4000\cdot10^{-24}+10^{-22}\bigr)\,\dd s
 <2\cdot10^{-14}.
\]
Thus the transverse component cannot trigger the exit in
\eqref{app:spiral:stopping-time}.

Finally, let
\(\mathcal R_K(y):=-e_1^\top S^{-1}b_K(z)\) be the point radial field.
Differentiating the same explicit \(2\times2\) inverse gives
\[
 \sup_{\norm{y}\leq10^{-12}}\norm{\nabla\mathcal R_K(y)}
 \leq\frac{20}{K},
 \qquad
 \mathcal R_K(0)=-\frac3{K+1}.
\]
The Gaussian correction is at most \(\varepsilon_b(t)\), so
\[
 \left|\frac{\dd}{\dd t}\log r_t+\frac3{K+1}\right|
 \leq\frac{20}{K}\norm{y_t}+\varepsilon_b(t)
 <\frac1{4K}.
\]
This proves \(-4/K\leq\dot r_t/r_t\leq-2/K\).  The radius can neither
reach the upper barrier nor, before time \(K/100\), fall below
\(e^{2-0.04}r_0>e^{3/2}r_0\).  Hence \(\tau=K/100\), which proves
\eqref{app:spiral:radial-shadowing}; integration gives
\eqref{app:spiral:mean-lower}.
\end{proof}

\subsubsection{Optimizer comparison and completion of the lower bound}

\begin{lemma}[Optimizer covariance and intrinsic condition number]
\label{app:spiral:optimizer-comparison}
Let \(C_\star\) be the optimizing covariance for \(V_K\), and define
\[
    \eta_K:=K\exp\left(-\frac{\alpha r_0^2}{2}\right).
\]
Under \eqref{app:spiral:radius-condition}, \(\eta_K\leq10^{-3}\) and
\begin{equation}
\label{app:spiral:optimizer-precision}
    \alpha_0(1-\eta_K)\Id
    \preceq C_\star^{-1}
    \preceq \alpha_0(1+\eta_K)\Id.
\end{equation}
Moreover,
\begin{equation}
\label{app:spiral:intrinsic-condition}
    K\leq\kappastar(V_K)\leq1.03K,
\end{equation}
and
\begin{equation}
\label{app:spiral:initial-product}
    \betastar\lambda_{0,\star}
    \geq\frac{1-\eta_K}{1+\eta_K}.
\end{equation}
\end{lemma}

\begin{proof}
The potential is exactly quadratic with Hessian \(\alpha_0\Id\) on
\(\{\norm{x}\leq r_0\}\).  Since
\(C_\star\preceq\alpha^{-1}\Id\), write
\(X_\star=C_\star^{1/2}Z\), \(Z\sim\N(0,\Id_2)\).  Then
\(\norm{X_\star}>r_0\) implies
\(\norm{Z}>\sqrt\alpha r_0\), and
\(\Prob(\norm{Z}>s)=e^{-s^2/2}\) in two dimensions.  Combining this tail
identity with
\(\norm{\nabla^2V_K-\alpha_0\Id}_{\op}\leq \alpha_0K\) and
\[
    C_\star^{-1}=\E[\nabla^2V_K(X_\star)]
\]
proves \eqref{app:spiral:optimizer-precision}.

The global Hessian certificate and
\eqref{app:spiral:optimizer-precision} give
\[
    \alphastar\geq\frac{0.98}{1+\eta_K},
    \qquad
    \betastar\leq\frac{K+0.02}{1-\eta_K},
\]
and hence \(\kappastar\leq1.03K\).  At the origin the Hessian is
\(\alpha_0\Id\), while at a valley point it has eigenvalue \(\alpha_0K\).  Therefore
\[
    \alphastar\leq \alpha_0\lambda_{\min}(C_\star),
    \qquad
    \betastar\geq \alpha_0K\lambda_{\min}(C_\star),
\]
which proves \(\kappastar\geq K\).

Finally,
\eqref{app:spiral:initial-covariance-band} gives
\(C_0\succeq(\alpha_0K)^{-1}\Id\), whereas
\(C_\star^{-1}\succeq \alpha_0(1-\eta_K)\Id\).  Thus
\(\lambda_{0,\star}\geq(1-\eta_K)/K\).  The valley eigenvalue and
\eqref{app:spiral:optimizer-precision} also give
\(\betastar\geq K/(1+\eta_K)\), proving
\eqref{app:spiral:initial-product}.
\end{proof}

\begin{proof}[Proof of \Cref{thm:global:spiral-lower}]
\Cref{app:spiral:hessian-certificate} proves the global \(C^3\)
regularity and curvature bounds, while
\eqref{app:spiral:initial-covariance-band} proves admissibility of the
anisotropic initialization.  The potential is even, so its optimizing
mean is \(m_\star=0\).  Strong convexity in the mean and global
optimality of \((0,C_\star)\) imply
\[
    \Delta(m,C)\geq\frac\alpha2\norm{m}^2
    \qquad\text{for every }(m,C).
\]
By \Cref{app:spiral:shadowing}, for \(0\leq t\leq K/100\),
\[
    \Delta(a_t)
    \geq\frac\alpha2e^{-0.08}\norm{m_0}^2.
\]
Therefore no level
\(\delta_{\loc}<(\alpha/2)e^{-0.08}\norm{m_0}^2\) can be reached before
\(K/100\).  The intrinsic comparison
\(\kappastar\leq1.03K\) in
\Cref{app:spiral:optimizer-comparison} gives
\[
    T_{\loc}\geq\frac K{100}
    \geq\frac{\kappastar(V_K)}{103}.
\]
\end{proof}

For \(d>2\), define
\[
    V_K^{(d)}(x)
    :=V_K(x_1,x_2)+\frac{\alpha_0}{2}\sum_{j=3}^d x_j^2
\]
and initialize the added coordinates at equilibrium.  The block-diagonal
subsystem is invariant.  Hadamard's determinant inequality shows that a
cross-covariance block can only decrease the Gaussian determinant while
leaving the separated potential expectation unchanged, so the optimizing
covariance is block diagonal as well.  The two-dimensional slow subsystem and
its intrinsic condition comparison are unchanged.

The construction can also meet any prescribed positive Hessian-Lipschitz
constant.  Differentiating
\eqref{app:spiral:hessian-decomposition}, using
\(\norm{\chi'''}_\infty\leq52.5\) and
\(\norm{g_K'}_\infty\leq4.376/\delta_K\), gives
\begin{equation}
\label{app:spiral:global-hessian-lipschitz}
    \operatorname{Lip}_{\op}(\nabla^2V_K)
    \leq10^8\frac{\alpha_0K^{5/2}}{r_0}.
\end{equation}
Thus, for \(L_H>0\), it is enough to choose
\[
    r_0\geq
    \max\left\{
        \frac{10^8\alpha_0K^{5/2}}{L_H},
        \frac{10^{25}K^{3/2}\sqrt{\log(eK)}}{\sqrt\alpha}
    \right\}.
\]
The initial radius is allowed to grow with \(K\), and with \(L_H^{-1}\)
in this version.  At \(L_H=0\) the target is quadratic, and this
construction does not assert a linear lower bound under a radius budget
fixed independently of \(L_H\).

\subsection{The one-dimensional bump train}
\label{app:bump}

This subsection proves \Cref{thm:global:bump-lower}.  Normalize
\(\alpha=1\), \(\beta=\kappa\), and write a scalar Gaussian parameter as
\(a=(m,c)\).  For \(X\sim\N(m,c)\), set
\begin{equation}
\label{app:bump:averaged-score-curvature}
    b(m,c):=\E[V'(X)],
    \qquad
    A(m,c):=\E[V''(X)].
\end{equation}
The two deterministic maps reduce respectively to
\begin{align}
\label{app:bump:riemannian-map}
    m^+&=m-hc\,b(m,c),
    &
    c^+&=c\exp\bigl(h[1-cA(m,c)]\bigr),\\
\label{app:bump:kl-map}
    m^+&=m-hc\,b(m,c),
    &
    c^+&=\frac{(1+h)c}{1+hcA(m,c)}.
\end{align}
Fix
\begin{equation}
\label{app:bump:stepsize}
    h=\frac\gamma\kappa,
    \qquad
    0<\gamma\leq\frac12,
\end{equation}
with \(\gamma\) independent of \(\kappa\).

\subsubsection{Bump profile, placement, and regularity}

Choose an even \(\varphi\in C_c^\infty(\R)\) such that
\begin{equation}
\label{app:bump:profile}
    0\leq\varphi\leq1,
    \qquad
    \varphi(u)=1\quad(|u|\leq\tfrac12),
    \qquad
    \varphi(u)=0\quad(|u|\geq1).
\end{equation}
Let
\[
    I_\varphi:=\int_\R\varphi(u)\,\dd u,
    \qquad
    M_\varphi:=\norm{\varphi'}_\infty.
\]
For a prescribed \(L_0>0\), define
\begin{equation}
\label{app:bump:width-mass}
    w_\kappa:=\frac{(\kappa-1)M_\varphi}{L_0},
    \qquad
    B_\kappa:=w_\kappa I_\varphi,
    \qquad
    H_\kappa:=(\kappa-1)B_\kappa.
\end{equation}
Then \(w_\kappa=\Theta(\kappa)\),
\(H_\kappa=\Theta(\kappa^2)\), and
\begin{equation}
\label{app:bump:lipschitz-normalization}
    \frac{\kappa-1}{w_\kappa}\norm{\varphi'}_\infty=L_0.
\end{equation}

Choose \(T>0\) sufficiently small and \(Y>0\) sufficiently large,
depending only on \(\varphi\) and \(L_0\), and set
\begin{equation}
\label{app:bump:count-initial-center}
    N_\kappa:=\left\lfloor\frac{T\kappa^2}{\gamma}\right\rfloor,
    \qquad
    x_0:=\frac{Y\kappa^4}{\gamma}.
\end{equation}
Define the positive centers recursively by
\begin{equation}
\label{app:bump:center-recursion}
    x_{j+1}
    :=x_j-\frac{\gamma}{\kappa^2}
    \left[
        x_j+H_\kappa\left(N_\kappa-j+\frac12\right)
    \right],
    \qquad 0\leq j<N_\kappa.
\end{equation}

\begin{lemma}[Geometry of the centers]
\label{app:bump:center-geometry}
For the preceding choices and all sufficiently large \(\kappa\),
\begin{equation}
\label{app:bump:center-lower}
    x_{N_\kappa}\geq\frac34x_0
\end{equation}
and
\begin{equation}
\label{app:bump:center-spacing}
    x_j-x_{j+1}\geq\frac{3Y}{4}\kappa^2,
    \qquad 0\leq j<N_\kappa.
\end{equation}
Hence the intervals
\([x_j-w_\kappa,x_j+w_\kappa]\) are pairwise disjoint and disjoint
from their reflections about the origin.
\end{lemma}

\begin{proof}
Summing \eqref{app:bump:center-recursion} and using \(x_j\leq x_0\)
gives
\begin{align*}
    x_0-x_{N_\kappa}
    &\leq
    \frac{\gamma N_\kappa}{\kappa^2}x_0
    +\frac{\gamma H_\kappa}{\kappa^2}
      \sum_{j=0}^{N_\kappa-1}
      \left(N_\kappa-j+\frac12\right)\\
    &\leq
    Tx_0+C_\varphi\frac{T^2\kappa^4}{\gamma},
\end{align*}
where \(C_\varphi\) depends only on \(\varphi,L_0\).  Since
\(x_0=Y\kappa^4/\gamma\), choosing \(T\) small and then \(Y\) large
makes the right side at most \(x_0/4\), proving
\eqref{app:bump:center-lower}.  The recursion and this lower bound give
\[
    x_j-x_{j+1}
    \geq\frac{\gamma}{\kappa^2}x_{N_\kappa}
    \geq\frac{3Y}{4}\kappa^2.
\]
Because \(w_\kappa=O(\kappa)\), the positive supports are disjoint for
large \(\kappa\).  Also
\(x_{N_\kappa}\asymp\kappa^4/\gamma\gg w_\kappa\), so they are disjoint
from the reflected supports.
\end{proof}

Define \(V_\kappa\) by
\begin{equation}
\label{app:bump:potential}
    V_\kappa''(x)
    :=1+(\kappa-1)
    \sum_{j=0}^{N_\kappa}
    \left[
        \varphi\left(\frac{x-x_j}{w_\kappa}\right)
        +\varphi\left(\frac{x+x_j}{w_\kappa}\right)
    \right],
\end{equation}
with \(V_\kappa'(0)=V_\kappa(0)=0\).

\begin{lemma}[Regularity and exact condition number]
\label{app:bump:regularity}
For all sufficiently large \(\kappa\), the potential \(V_\kappa\) is
even and belongs to \(C^\infty(\R)\).  Moreover,
\begin{equation}
\label{app:bump:curvature-bounds}
    1\leq V_\kappa''(x)\leq\kappa
    \qquad(x\in\R),
\end{equation}
both endpoints are attained, and
\begin{equation}
\label{app:bump:hessian-lipschitz}
    \norm{V_\kappa'''}_\infty\leq L_0.
\end{equation}
The density \(e^{-V_\kappa}\) is integrable, and the original and
optimizer-whitened condition numbers are both \(\kappa\).
\end{lemma}

\begin{proof}
The reflected construction makes \(V_\kappa''\) even, so the
normalizations at zero make \(V_\kappa'\) odd and \(V_\kappa\) even.
By \Cref{app:bump:center-geometry}, at most one summand in
\eqref{app:bump:potential} is nonzero at any point.  Hence
\(1\leq V_\kappa''\leq1+(\kappa-1)=\kappa\).  The lower value is
attained outside the supports, and the upper value on each flat top.
Likewise,
\[
    |V_\kappa'''(x)|
    \leq\frac{\kappa-1}{w_\kappa}\norm{\varphi'}_\infty=L_0.
\]
Since \(V_\kappa''\geq1\) and
\(V_\kappa(0)=V_\kappa'(0)=0\), one has
\(V_\kappa(x)\geq x^2/2\), proving integrability.  In one dimension,
optimizer whitening multiplies every Hessian value by the same positive
scalar \(c_\star\), so it does not change the ratio of the largest and
smallest Hessian values.  Thus \(\kappastar=\kappa\).
\end{proof}

\subsubsection{Gaussian averages on the train}

Set
\begin{equation}
\label{app:bump:matched-covariance}
    c_\kappa:=\frac1\kappa
\end{equation}
and define ideal states \(a_j^\circ=(x_j,c_\kappa)\).

\begin{lemma}[Averaged score and curvature]
\label{app:bump:averages}
There exist constants \(C,c>0\) and an integer \(p\geq1\), independent
of \(\kappa\), such that, for \(0\leq j\leq N_\kappa\),
\begin{align}
\label{app:bump:center-curvature}
    A(x_j,c_\kappa)
    &=\kappa+r^A_{j,\kappa},
    &
    |r^A_{j,\kappa}|
    &\leq C\kappa^p e^{-c\kappa^3},\\
\label{app:bump:center-score}
    b(x_j,c_\kappa)
    &=x_j+H_\kappa\left(N_\kappa-j+\frac12\right)
      +r^b_{j,\kappa},
    &
    |r^b_{j,\kappa}|
    &\leq C\kappa^p e^{-c\kappa^3}.
\end{align}
The same estimates are uniform on
\begin{equation}
\label{app:bump:tube}
    \mathcal T_{j,\kappa}
    :=\left\{
        (m,c):
        |m-x_j|\leq\frac{w_\kappa}{8},
        \quad
        |\kappa c-1|\leq\frac18
    \right\}.
\end{equation}
More precisely, for \((m,c)\in\mathcal T_{j,\kappa}\),
\begin{align}
\label{app:bump:tube-curvature}
    |A(m,c)-\kappa|
    &\leq\varepsilon_\kappa,\\
\label{app:bump:tube-score}
    \left|
        b(m,c)-x_j
        -H_\kappa\left(N_\kappa-j+\frac12\right)
        -\kappa(m-x_j)
    \right|
    &\leq\varepsilon_\kappa,
\end{align}
where
\begin{equation}
\label{app:bump:epsilon}
    \varepsilon_\kappa:=C\kappa^p e^{-c\kappa^3}.
\end{equation}
\end{lemma}

\begin{proof}
At the ideal state write \(X=x_j+Z/\sqrt\kappa\), with
\(Z\sim\N(0,1)\).  The current bump equals one whenever
\(|X-x_j|\leq w_\kappa/2\), and \(w_\kappa\asymp\kappa\), so
\[
    \Prob\left(|X-x_j|>\frac{w_\kappa}{2}\right)
    \leq2e^{-c\kappa^3}.
\]
Every other bump is at distance at least \(c\kappa^2\), or
\(c\kappa^{5/2}\) Gaussian standard deviations, from \(x_j\).
Summing over \(N_\kappa+1=O(\kappa^2)\) bumps changes only the
polynomial prefactor.  This proves
\eqref{app:bump:center-curvature}.

Let
\[
    \Phi(s):=\int_{-\infty}^s\varphi(u)\,\dd u.
\]
For \(x>0\), except on an exponentially unlikely event on which the
Gaussian sample crosses the origin, integration of
\eqref{app:bump:potential} gives
\[
    V_\kappa'(x)
    =x+(\kappa-1)w_\kappa
      \sum_{k=0}^{N_\kappa}
      \Phi\left(\frac{x-x_k}{w_\kappa}\right).
\]
The bumps with \(k>j\) lie between the origin and \(x_j\) and contribute
their full mass \(I_\varphi\); those with \(k<j\) contribute zero, up to
the same Gaussian-tail error.  For the current bump,
\(\Phi(s)+\Phi(-s)=I_\varphi\), and symmetry gives
\[
    \E\Phi\left(\frac{X-x_j}{w_\kappa}\right)
    =\frac{I_\varphi}{2}.
\]
Multiplication by \((\kappa-1)w_\kappa\) proves
\eqref{app:bump:center-score}.

Now take \((m,c)\in\mathcal T_{j,\kappa}\).  Its mean stays at least
\(3w_\kappa/8\) from the edge of the current flat top, and its standard
deviation is at most \(\sqrt{9/(8\kappa)}\).  The probability of
reaching a transition layer remains bounded by \(Ce^{-c\kappa^3}\),
uniformly on the tube.  This proves
\eqref{app:bump:tube-curvature}.  Since
\(\partial_m b(m,c)=A(m,c)\), integration in \(m\) gives the term
\(\kappa(m-x_j)\) with error \(\varepsilon_\kappa\).  Price's identity
also gives
\[
    \partial_c b(m,c)=\frac12\E[V_\kappa'''(X)].
\]
The third derivative is supported on transition layers, which are
reached with probability \(O(e^{-c\kappa^3})\) throughout the tube.
Integration in \(c\) proves \eqref{app:bump:tube-score}, after enlarging
the polynomial prefactor in \eqref{app:bump:epsilon}.
\end{proof}

\subsubsection{Uniform shadowing for both covariance maps}

Let \((m_j,c_j)\) be the iterates of either
\eqref{app:bump:riemannian-map} or \eqref{app:bump:kl-map}, initialized
by
\begin{equation}
\label{app:bump:initial-state}
    m_0=x_0,
    \qquad
    c_0=c_\kappa=\frac1\kappa.
\end{equation}
Define
\begin{equation}
\label{app:bump:shadow-errors}
    e_j:=m_j-x_j,
    \qquad
    q_j:=\kappa c_j-1.
\end{equation}

\begin{lemma}[Uniform discrete shadowing]
\label{app:bump:shadowing}
For either covariance map, there exist constants \(C,c>0\) and an
integer \(p\geq1\) such that, for all sufficiently large \(\kappa\) and
every \(0\leq j\leq N_\kappa\),
\begin{equation}
\label{app:bump:shadowing-estimate}
    |e_j|+|q_j|\leq C\kappa^p e^{-c\kappa^3}.
\end{equation}
In particular, every iterate lies in \(\mathcal T_{j,\kappa}\) and
\[
    m_j=x_j+o(1),
    \qquad
    c_j=\frac1\kappa(1+o(1))
\]
uniformly for \(0\leq j\leq N_\kappa\).
\end{lemma}

\begin{proof}
Write
\[
    B_j:=x_j+H_\kappa\left(N_\kappa-j+\frac12\right).
\]
The center recursion says
\(x_{j+1}=x_j-\gamma B_j/\kappa^2\).  On the tube,
\Cref{app:bump:averages} gives
\[
    b(m_j,c_j)=B_j+\kappa e_j+r_j,
    \qquad
    |r_j|\leq\varepsilon_\kappa.
\]
Since \(c_j=(1+q_j)/\kappa\), the mean error obeys the exact recursion
\begin{equation}
\label{app:bump:mean-error-recursion}
    e_{j+1}
    =\left(1-\frac{\gamma(1+q_j)}{\kappa}\right)e_j
      -\frac{\gamma q_j}{\kappa^2}B_j
      -\frac{\gamma(1+q_j)}{\kappa^2}r_j.
\end{equation}
The center geometry gives \(|B_j|\leq C\kappa^4/\gamma\).  As long as
\(|q_j|\leq1/8\),
\begin{equation}
\label{app:bump:mean-error-bound}
    |e_{j+1}|
    \leq
    \left(1-\frac{7\gamma}{8\kappa}\right)|e_j|
    +C\kappa^2|q_j|
    +C\frac{\varepsilon_\kappa}{\kappa^2}.
\end{equation}

For the Riemannian covariance map write
\(A(m_j,c_j)=\kappa+\delta_j\), with
\(|\delta_j|\leq\varepsilon_\kappa\).  Then
\begin{equation}
\label{app:bump:q-riemannian-exact}
    1+q_{j+1}
    =(1+q_j)
      \exp\left[
        \frac{\gamma}{\kappa}
        \left(
            1-(1+q_j)\left(1+\frac{\delta_j}{\kappa}\right)
        \right)
      \right].
\end{equation}
For \(|q_j|\leq1/8\) and large \(\kappa\), the derivative of the right
side with respect to \(q_j\) lies between zero and
\(1-\gamma/(2\kappa)\), while its variation with respect to
\(\delta_j\) is at most \(C\gamma|\delta_j|/\kappa^2\).  Since the map
fixes \(q=0\) when \(\delta=0\), the mean-value theorem yields
\begin{equation}
\label{app:bump:q-riemannian-bound}
    |q_{j+1}|
    \leq\left(1-\frac{\gamma}{2\kappa}\right)|q_j|
    +C\frac{\varepsilon_\kappa}{\kappa^2}.
\end{equation}

For the KL/Bregman covariance map,
\begin{equation}
\label{app:bump:q-kl-exact}
    1+q_{j+1}
    =
    \frac{(1+\gamma/\kappa)(1+q_j)}
         {1+(\gamma/\kappa)(1+q_j)(1+\delta_j/\kappa)}.
\end{equation}
The same mean-value argument gives
\begin{equation}
\label{app:bump:q-kl-bound}
    |q_{j+1}|
    \leq\left(1-\frac{\gamma}{2\kappa}\right)|q_j|
    +C\frac{\varepsilon_\kappa}{\kappa^2}.
\end{equation}

Since \(q_0=0\), either covariance recursion implies
\begin{equation}
\label{app:bump:q-summed}
    |q_j|
    \leq C\frac{\varepsilon_\kappa}{\gamma\kappa}
    \qquad(0\leq j\leq N_\kappa).
\end{equation}
Substitute this into \eqref{app:bump:mean-error-bound}, use \(e_0=0\),
and sum the geometric recursion to obtain
\begin{equation}
\label{app:bump:e-summed}
    |e_j|\leq C_\gamma\kappa^p e^{-c\kappa^3}
\end{equation}
after increasing \(p\) if necessary.  Exponential decay dominates
every polynomial, so
\eqref{app:bump:q-summed}--\eqref{app:bump:e-summed} are eventually
smaller than \(1/8\) and \(w_\kappa/8\).  The tube hypotheses close by
induction, proving \eqref{app:bump:shadowing-estimate}.
\end{proof}

\subsubsection{Objective comparison and completion}

For \(c>0\), define
\begin{equation}
\label{app:bump:smoothed-potential}
    F_c(m):=\E_{X\sim\N(m,c)}[V_\kappa(X)].
\end{equation}
Up to a target-dependent additive constant,
\begin{equation}
\label{app:bump:scalar-objective}
    \calE_\kappa(m,c)=F_c(m)-\frac12\log c.
\end{equation}

\begin{lemma}[Constant-factor objective lower bound]
\label{app:bump:objective-gap}
There is a numerical \(q\in(0,1)\), independent of \(\kappa\), such
that the iterates of either scheme satisfy
\begin{equation}
\label{app:bump:gap-lower}
    \Delta(m_j,c_j)\geq q\,\Delta(m_0,c_0),
    \qquad 0\leq j\leq N_\kappa,
\end{equation}
for all sufficiently large \(\kappa\).
\end{lemma}

\begin{proof}
The potential is even, so \(F_c\) is even and \(F_c'(0)=0\).  Moreover,
\(F_c''(m)=A(m,c)\geq1\), whence
\begin{equation}
\label{app:bump:smoothed-lower}
    F_c(m)-F_c(0)\geq\frac12m^2.
\end{equation}
At the matched covariance,
\Cref{app:bump:center-geometry} gives
\begin{equation}
\label{app:bump:final-mean-lower}
    F_{c_\kappa}(x_j)-F_{c_\kappa}(0)
    \geq\frac12x_j^2
    \geq\frac12x_{N_\kappa}^2
    \geq\frac9{32}x_0^2.
\end{equation}

Integration of \eqref{app:bump:potential} gives the global estimate
\begin{equation}
\label{app:bump:global-score-bound}
    |V_\kappa'(x)|
    \leq |x|+H_\kappa(N_\kappa+1).
\end{equation}
Therefore, for \(m\geq0\),
\[
    F_{c_\kappa}'(m)
    \leq
    m+\sqrt{\frac{2}{\pi\kappa}}
      +H_\kappa(N_\kappa+1).
\]
Integrating from zero to \(x_0\) gives
\begin{equation}
\label{app:bump:initial-mean-upper}
    F_{c_\kappa}(x_0)-F_{c_\kappa}(0)
    \leq
    \frac12x_0^2
    +\left[
        H_\kappa(N_\kappa+1)
        +\sqrt{\frac{2}{\pi\kappa}}
    \right]x_0
    \leq C_0x_0^2,
\end{equation}
where \(C_0\) is independent of \(\kappa\); here
\(H_\kappa(N_\kappa+1)=O(\kappa^4/\gamma)\) and
\(x_0=Y\kappa^4/\gamma\).  Combining
\eqref{app:bump:final-mean-lower} and
\eqref{app:bump:initial-mean-upper} yields
\begin{equation}
\label{app:bump:mean-ratio}
    F_{c_\kappa}(x_j)-F_{c_\kappa}(0)
    \geq q_0
    [F_{c_\kappa}(x_0)-F_{c_\kappa}(0)],
    \qquad
    q_0:=\frac9{32C_0}>0.
\end{equation}

For every fixed covariance,
\(m\mapsto\calE_\kappa(m,c)\) is strictly convex and even, so its
unique minimizing mean is zero.  Let \(c_\star\) be the optimizing
covariance and set
\[
    D_\kappa
    :=\calE_\kappa(0,c_\kappa)-\calE_\kappa(0,c_\star)\geq0.
\]
At covariance \(c_\kappa\),
\[
    \Delta(x_j,c_\kappa)
    =F_{c_\kappa}(x_j)-F_{c_\kappa}(0)+D_\kappa.
\]
After replacing \(q_0\) by \(\min\{q_0,1\}\),
\eqref{app:bump:mean-ratio} implies
\begin{equation}
\label{app:bump:ideal-gap-ratio}
    \Delta(x_j,c_\kappa)\geq q_0\Delta(x_0,c_\kappa).
\end{equation}

It remains to replace the ideal states by the actual iterates.  On the
tube,
\[
    |\partial_m\calE_\kappa|=|b(m,c)|
    \leq C\kappa^4/\gamma,
\]
while
\[
    \left|\partial_c\calE_\kappa(m,c)\right|
    =\frac12\left|A(m,c)-\frac1c\right|
    \leq C(\varepsilon_\kappa+\kappa|q|).
\]
Together with \Cref{app:bump:shadowing}, the mean-value theorem gives
\begin{equation}
\label{app:bump:objective-shadow}
    \sup_{0\leq j\leq N_\kappa}
    \left|
        \calE_\kappa(m_j,c_j)
        -\calE_\kappa(x_j,c_\kappa)
    \right|
    =o(x_0^2).
\end{equation}
By \eqref{app:bump:smoothed-lower},
\(\Delta(x_0,c_\kappa)\geq x_0^2/2\), so the error in
\eqref{app:bump:objective-shadow} is
\(o(\Delta(x_0,c_\kappa))\).  Absorbing it in
\eqref{app:bump:ideal-gap-ratio}, after reducing \(q_0\) by a fixed
factor, proves \eqref{app:bump:gap-lower}.
\end{proof}

\begin{proof}[Proof of \Cref{thm:global:bump-lower}]
\Cref{app:bump:regularity} proves
\(1\leq V_\kappa''\leq\kappa\),
\(\norm{V_\kappa'''}_\infty\leq L_0\), and
\(\kappastar=\kappa\).  The initialization
\eqref{app:bump:initial-state} lies at the lower edge of the
curvature-scale band \([1/\kappa,1]\), so the construction begins after
covariance burn-in.  By \Cref{app:bump:objective-gap},
\[
    \Delta_n\geq q\Delta_0
    \qquad
    \left(0\leq n\leq
    N_\kappa=\left\lfloor\frac{T\kappa^2}{\gamma}\right\rfloor\right)
\]
for either exact-expectation map.  Since \(h=\gamma/\kappa\),
\[
    \frac{\kappa^2}{\gamma}=\frac{\kappa}{h},
\]
which proves the claimed \(\Omega(\kappa/h)\) constant-factor lower
bound for both discretizations.
\end{proof}

The construction is tied to the fixed value of \(\gamma\) and does not
cover adaptive relative-curvature steps, line searches, or implicit updates.
Indeed, along the shadowed path,
\[
    c_jA(m_j,c_j)=1+o(1),
    \qquad
    c_j\simeq\frac1\kappa,
    \qquad
    \frac{b(m_j,c_j)}{m_j}=\Theta(1).
\]
Thus the relative mean motion is
\[
    h\frac{c_jb(m_j,c_j)}{m_j}
    =\Theta\left(\frac{\gamma}{\kappa^2}\right),
\]
which is the precise source of the fixed-step slowdown.
